Krátké, ověřené příklady z praxe slouží jako inspirace pro pilotní nasazení na VUT; níže jsou tři konkrétní příklady‑vzory s praktickými kroky a doporučením, co lze adaptovat lokálně. Některé širší příklady institucionálních nasazení (např. Georgia Tech, Harvard CS50, Arizona State) jsou užitečnou inspirací, jejich popisy zde uvádíme jako obecné pozadí (general background knowledge).

\medskip

\textbf{1) Minecraft Education — výuka principů AI hrou.}  
Minecraft Education nabízí připravené světy a aktivity vhodné pro demonstraci principů strojového učení a pro „Hour of Code“ aktivity; některé z těchto hodin jsou dostupné zdarma, další jsou součástí školních licencí Microsoft 365 A3/A5 \cite{kb_ai_101_tipu_pdf_04e4a115_page_15_2}.  
Praktický pilot (s návrhem kroků):
\begin{enumerate}
  \item Stažení a výběr vhodného světa v Minecraft Education (např. hodina zaměřená na AI nebo Hour of Code) \cite{kb_ai_101_tipu_pdf_04e4a115_page_15_2}.
  \item Příprava krátkého doprovodného materiálu pro studenty: základní pojmy (co je model, trénink, data) a očekávané výstupy z aktivity (diskuse, krátká reflexe).
  \item Vedení hodiny: praktická aktivita ve světě, následná diskuse o možnostech a omezeních AI a závěrečné shrnutí.
\end{enumerate}
Krátké doporučení pro VUT (general background knowledge): vyzkoušet aktivitu v dílčím semináři nebo laboratoři a vyhodnotit pedagogický přínos před širším zavedením; zohlednit potřebu licence Microsoft 365 pro úplnou sadu obsahů.

\medskip

\textbf{2) Automatická tvorba testů: Microsoft 365 Copilot + Forms.}  
Microsoft 365 Copilot ve spojení s Forms umí vygenerovat kompletní test (otázky typu ABC, pravda/nepravda, otevřené), včetně označení správných odpovědí a vysvětlení, proč je odpověď správná \cite{kb_ai_101_tipu_pdf_04e4a115_page_36_1}.  
Praktický pilot (s návrhem kroků):
\begin{enumerate}
  \item Definujte učební cíl a rozsah testu (např. konkrétní kapitola nebo tematický blok).
  \item Vytvořte prompt pro Copilot: počet otázek, typy položek, úroveň (studenti 1.–2. ročníku apod.) a požadavek na klíč a krátké vysvětlení správných odpovědí \cite{kb_ai_101_tipu_pdf_04e4a115_page_36_1}.
  \item Zkontrolujte a upravte vygenerované otázky ručně (ověření faktické přesnosti a vhodnosti pro hodnocení).
  \item Publikujte přes Forms a po odevzdání využijte automatické značení tam, kde je to možné.
\end{enumerate}
Praktické poznámky pro VUT (general background knowledge): Copilot‑workflow zrychlí přípravu písemek a banky úloh; vždy však požadujte lidskou finalizaci otázek a kontrolu vhodnosti pro daný kurs.

\medskip

\textbf{3) Podpora procvičování v matematice — „Pokrok v matematice“ a OneNote.}  
Funkce „Pokrok v matematice“ umožňuje pedagogovi generovat příklady z vybrané oblasti, sbírat odpovědi a (volitelně) postupy žáků; nástroj poskytuje předběžné opravy a diagnostiku oblastí, kde studenti chybují. OneNote má podporu převodu rukopisu na digitální rovnice včetně výpočtu a zobrazení postupu \cite{kb_ai_101_tipu_pdf_04e4a115_page_8_1}.  
Praktický pilot (s návrhem kroků):
\begin{enumerate}
  \item Vyberte téma (např. lineární rovnice) a nastavte v nástroji generování s požadavkem na nahrání postupu řešení od studentů \cite{kb_ai_101_tipu_pdf_04e4a115_page_8_1}.
  \item Pilotně povolte nahrávání fotografií postupu nebo použití OneNote pro převod rukopisu; využijte automatické předopravování jako pomoc učiteli při předkontrole \cite{kb_ai_101_tipu_pdf_04e4a115_page_8_1}.
  \item Definujte pravidla eskalace (kdy poslat práci k lidské revizi) a provádějte spot‑checky kvality výsledků.
\end{enumerate}
Doporučení pro VUT (general background knowledge): integrace do seminářů a cvičení v laboratořích s technickým zaměřením; zaměřit se na využití výstupů AI především jako podpory formativní zpětné vazby.

\medskip

Závěrem: výše uvedené příklady ukazují praktické, nízkonákladové piloty, které lze replikovat na úrovni předmětů nebo laboratoří. Inspirace z větších institucionálních případů (CS50, campus‑licence ChatGPT Enterprise apod.) je užitečná pro strategické plánování, nicméně konkrétní kroky a smluvní ujednání je nutné konzultovat s interními IT a GDPR týmy na VUT (general background knowledge).